\documentclass[10pt,a4paper]{amsart}
\usepackage[margin=19mm]{geometry}
\usepackage{amsmath,amssymb,mathtools}
\usepackage[scr=rsfs,cal=euler]{mathalfa}
\usepackage{xfrac}
\usepackage[unicode,hidelinks,bookmarks=false]{hyperref}
\newcommand{\ie}{i.e.\ }
\newcommand{\cf}{cf.\ }
\newcommand{\resp}{respectively\ }
\newcommand{\Subsection}{\subsection}
\providecommand{\nobreakdash}{\nobreak\hbox{-}}
\setlength{\emergencystretch}{2em}
\multlinegap0pt
\usepackage{bm}
\usepackage{bbm}
\usepackage{dsfont}
\usepackage{xspace}
\usepackage{cancel}
\usepackage{mathtools}
\usepackage{amsthm}
\makeatletter
\@ifundefined{theorem}{\newtheorem{theorem}{Theorem}}{}
\@ifundefined{lemma}{\newtheorem{lemma}[theorem]{Lemma}}{}
\makeatother
\title[Testing properties of trees in graphical models with covaria]{Testing properties of trees in graphical models with covariance queries}
\author{Sofiya Burova, Francisco Calvillo, G\'abor Lugosi, Piotr Zwiernik}
\date{Source: arXiv:2605.15996v1 (CC BY 4.0)}
\begin{document}
\maketitle

\noindent\emph{Source and licence.} Testing properties of trees in graphical models with covariance queries. Version arXiv:2605.15996v1 (CC BY 4.0), distributed under the Creative Commons Attribution 4.0 International licence, as recorded in the arXiv licence metadata for this exact version and shown on the abstract page https://arxiv.org/abs/2605.15996v1. Full rights notice: licenses/arxiv-2605.15996-CC-BY-4.0.txt.\par\smallskip

\noindent\textit{Editorial scope.} Counted unit: the Lemma whose source label is \texttt{lem:TX} in the article (source0.tex lines 655--657, source label \texttt{lem:TX}), delivered together with the article's own complete original proof of it (source0.tex lines 658--694, one \texttt{proof} environment of 37 lines). No mathematics is written, repaired or re-derived: statement and proof are the article's own text line for line. Required context retained uncounted: none. Every definition, display and label of those lines is retained unchanged. Omitted from this package and not redistributed: 1--654, 695--1189. Nothing inside a retained excerpt is omitted or altered: every retained body is delivered exactly as the article's lines are written, so no omission receipt of any kind accompanies this package. Citations: the retained excerpts carry no citation command at all, so no bibliography entry of the article is transcribed and no reference is redistributed. Reference adaptation: none was needed and none was made; every reference inside the delivered unit resolves inside it, so no unresolved reference or citation is produced. Printed slips kept literal: no printed word, symbol, hypothesis or number of the counted statement or of its proof was altered. Layout and class: the article's own class is \texttt{article} with packages including kpfonts, titlesec, geometry, dsfont, comment, textcomp, amssymb, bbm, fancyhdr, amsmath, amsbsy, amsthm, amscd, latexsym; the wrapper is the lane's pinned \texttt{amsart} editorial wrapper at 10pt on A4, which restates the theorem-like environments the retained text needs and adds the article's own macro definitions that the retained text needs (none), with extra wrapper packages (bm, bbm, dsfont, xspace, cancel, mathtools). The delivered numbering is the unit's own. Assets: no retained span contains a figure environment, a graphics-inclusion command, a PDF-page inclusion, a table or a TikZ picture; no image file is copied into this package, the package declares no asset dependency and no image file is redistributed with it. Licence: version arXiv:2605.15996v1 of this article is distributed under the Creative Commons Attribution 4.0 International licence, as recorded in the arXiv licence metadata for this exact version and shown on the abstract page https://arxiv.org/abs/2605.15996v1. A full rights notice is delivered at licenses/arxiv-2605.15996-CC-BY-4.0.txt. Characters: every retained excerpt is pure ASCII, so the delivered engine, XeLaTeX with its default Latin Modern text font, reports no missing character.
\begin{lemma}\label{lem:TX}
Let $T$ be a weighted tree on vertex set $V$ and let $X\subseteq V$. Given all distances $d(x,v)$ with $x\in X$ and $v\in V$, one can recover the minimal subtree $T_X$ of $T$ spanning $X$, together with all its edge lengths. Moreover, one can recover all distances $d(u,v)$ with $u\in V(T_X)$ and $v\in V$. The query complexity is $n|X|$.
\end{lemma}

 \begin{proof}
     We have $w\in V(T_X)$ if and only if $d(x,x')=d(x,w)+d(w,x')$ for some $x,x'\in X$. This identifies the vertex set $V(T_X)$ of $T_X$. For the edges, note that $uv\in E(T_X)$ if and only if 
     \begin{itemize}
         \item [(i)] $u,v\in \overline{xx'}$ for some $x,x'\in X$,
         \item [(ii)] if $w\in \overline{xx'}$, then $w$ is closer to $x$, or closer to $x'$ than both $u$ and $v$.
     \end{itemize}
     Both conditions easily translate in terms of the queried distances:
     \begin{itemize}
         \item [(i)] $d(x,x')=d(x,u)+d(u,x')=d(x,v)+d(v,x')$ for some $x,x'\in X$
         \item [(ii)] $d(x,x')=d(x,w)+d(w,x')$ implies that 
         $$d(x,w)\leq \min\{d(x,u),d(x,v)\} \quad \text{or} \quad d(w,x')\leq \min\{d(u,x'),d(v,x')\}~.$$
     \end{itemize}
     If $uv\in E(T_X)$ and $u,v\notin X$, we compute $d(u,v)$ without querying it using the formula $d(u,v)=|d(u,x)-d(v,x)|$. 

It remains to recover distances from $V(T_X)$ to vertices outside $V(T_X)$. Fix $v\in V\setminus V(T_X)$. There is a unique vertex $a(v)\in V(T_X)$ at which the path from $v$ to $V(T_X)$ first meets $V(T_X)$. We claim that $a(v)$ is the unique vertex $a\in V(T_X)$ such that $d(x,v)-d(x,a)$ is independent of $x\in X$. If $a=a(v)$, then every path from $x\in X$ to $v$ passes through $a$, and hence
$$
d(x,v)=d(x,a)+d(a,v)~.
$$
Thus $d(x,v)-d(x,a)$ is constantly equal to $d(a(v),v)$. Conversely, suppose $a\neq a(v)$. Let $e$ be the first edge on the path from $a$ to $a(v)$ in $T_X$. Removing $e$ separates $T_X$ into two components, both containing vertices of $X$ by minimality of $T_X$. Choose $x,y\in X$ on opposite sides of $e$, with $x$ on the side containing $a(v)$ and $y$ on the side containing $a$. For $y$, the path from $y$ to $v$ passes through $a$, so
$$
d(y,v)-d(y,a)=d(a,a(v))+d(a(v),v)~.
$$
For $x$, the path from $x$ to $v$ passes through $a(v)$, while the path from $x$ to $a$ also passes through $a(v)$, and hence
$$
d(x,v)-d(x,a)=d(a(v),v)-d(a(v),a)~.
$$
The two quantities differ by $2d(a,a(v))>0$. Hence $d(x,v)-d(x,a)$ is not constant over $X$. This proves the claim. Thus $a(v)$ is identifiable from the queried distances and the already recovered distances inside $T_X$. Once $a(v)$ is known, we compute
$$
d(a(v),v)=d(x,v)-d(x,a(v))
$$
for any $x\in X$; this is known even when $a(v)\notin X$, because $d(x,a(v))$ has already been recovered inside $T_X$. Finally, for every $u\in V(T_X)$,
$$
d(u,v)=d(u,a(v))+d(a(v),v)~,
$$
where $d(u,a(v))$ is known from the recovered subtree $T_X$. This recovers all distances $d(u,v)$ with $u\in V(T_X)$ and $v\in V$. The number of queried distances is exactly $n|X|$.
     % For the second part, if $u,v\in V(T_X)$ then $d(u,v)$ is identified from the first part. If $v\in V(T)\setminus V(T_X)$, let $x\in X$ be the closest to $v$. Then $d(u,v)$ can be computed by the formula $d(u,v)=d(u,x)+d(x,v)$.
 \end{proof}

\end{document}
